\documentclass{article}

\usepackage{../handouts}

\title{CSCI 535 Handout 01: Analyzing an Algorithm}
\author{Name: \rule{3in}{0.15mm}}
\date{3 September 2026}



\begin{document}
\maketitle

\noindent
{\bf What algorithm are you analyzing today?}

\textbf{Answer}
\vspace{1in}

\section{What?}

When analyzing algorithms, we first ask WHAT?  That is, what is the problem that
we want solved?  Typically, in this description, we need to understand both what
are the inputs \emph{and} what are the outputs of the algorithm.  When
describing WHAT, it should be independent of HOW and HOW FAST.

So, describe the WHAT for the algorithm you wish to analyze today.

What is the input?

\textbf{Answer}
\vspace{0.75in}

What is the output?

\textbf{Answer}
\vspace{0.75in}

Use the previous two answers to concisely describe what is the problem?  This
often requires explaining the output in terms of the input.

\textbf{Answer}
\vspace{1in}

\section{How?}

Often, it helps to walk through small examples.  If you came up with the
algorithm, this is a nice sanity check.  If you are studying an algorithm given
to you, this is helpful to better understand the algorithm.  Do that now.

\textbf{Answer}
\vspace{3in}

If you are presenting an algorithm, you should describe how it works in
pseudocode.  More importantly, explain what is going on in the pseudocode in the
prose text.  The prose should point to the algorithm.  (Remember, in LaTex,
algorithms are floating environments, so we need to be reminded to look at
them). If you are working on analyzing an algorithm that is given to you in
pseudocode, briefly explain HOW it works here.  In a first pass, ignore base cases/small
things and
focus on, \emph{what is the key step}? What really makes this algorithm work?
Then, go back and fill in the details.

\textbf{Answer}
\vspace{3in}


\pagebreak
\section{How Fast?}

\subsection{Option A: Asymptotics}
What is the asymptotic runtime?  If this is a recursive algorithm, you are expected to
explicitly state the recurrence relation.  Always use asymptotic notation, being
as precise as possible: don't use big-O notation when you can use big-Omega.

\textbf{Answer}
\vspace{1in}

\subsection{Option B: Decrementing Functions}
Sometimes, it is hard to pin down the runtime, yet we still want to ensure that
the algorithm terminates.  For this, we use decrementing functions to prove that
loops/recursions terminate.

A decrementing function is used to prove that an algorithm (specifically, a loop
or a recurrence) terminates.  We use
them in two ways: (1) when doing back-of-hand calculations to ensure that our
algorithm will terminate before doing a proper runtime analysis; (2) when it is
difficult to prove exact runtime.

\begin{definition}[Decrementing Function]
    A decrementing function is a function\footnote{In the decrementing function,
    $\mathcal{S}$ is the state space, and $\N$ denotes the natural numbers.  We
    can actually replace $\N$ with any well-ordered set though.}
    $D \colon \mathcal{S} \to \N$ that
    satisfies the following properties:
    \begin{enumerate}[(1)]
        \item Each time a recursive call is made, the function strictly
            decreases between the top of the parent call to the top of the child
            call.
        \item Each time a loop is re-entered, the function is strictly less than
            the last time it entered that loop.
    \end{enumerate}
\end{definition}

If such a function exists, then our algorithm terminates!  Can you explain why
this is true?


As an exercise, even if you know the runtime, try to use a
decrementing function to prove termination.

\begin{enumerate}
    \item What is the decrementing function for this algorithm?
    \practice

    \item Justify (=prove) why function is well-defined.  That is, why is the
        output of the function a natural number?
    \practice

    \item     Justify (=prove) why this decrements each time it reaches the top
        of the loop (or each time it goes into a new recursive call).
    \practice

\end{enumerate}


\pagebreak
\section{Why? The Loop Invariant}

Finally, we prove why the algorithm works.  Typically, this is done with a loop
(or recursion) invariant.  Here, we focus on setting up the loop invariant.
What
are the following statements (start with your best guess, then come back and
revise if needed):

The loop guard $G$:
\vspace{0.5in}

The post-condition $Q$:
\vspace{0.5in}

The pre-condition $P$:
\vspace{0.5in}

The loop invariant $L=L_i$:
\vspace{0.5in}


\pagebreak
Next, we use use this loop invariant (and other statements above)
to prove partial correctness There are three parts to partial
correctness: Initialization, Maintenance, and End.

Initialization: If the preconditions are met and if we enter the loop, then the
loop invariant is correct. ($P \wedge G \implies L$).

\textbf{Answer}
\vspace{1in}


Maintenance: If we entered the loop and the loop invariant held, then when we
exit the loop, the loop invariant will still hold: ($G \wedge L_i \implies
L_{i+1})$.

\textbf{Answer}
\vspace{1in}


End: If the loop ends and the loop invariant was correct, then the algorithm is
correct. ($L \wedge \neg G \implies Q$).

\textbf{Answer}
\vspace{1in}

\end{document}
